%% sections/conclusion.tex

\section{Conclusion}
\label{sec:conclusion}

This paper presented RADAR, a multimodal system for distinguishing human-authored texts from AI-generated content, submitted to the Voight-Kampff Generative AI Detection task at PAN 2026. RADAR couples a powerful pre-trained transformer encoder (ModernBERT-large) with a lightweight statistical linguistic profiler through a cross-attention fusion layer. The system is trained in two stages: an initial encoder-frozen phase that stabilises the lightweight components, followed by full end-to-end fine-tuning that allows the encoder to adapt to the detection objective.

Our main findings can be summarised as follows. First, the combination of deep contextual representations and handcrafted stylometric features is mutually beneficial: each stream captures aspects of AI authorship that the other misses, and cross-attention fusion allows the semantic stream to selectively weigh stylometric evidence. Second, a diverse training corpus is essential for generalisation: augmenting the organiser-provided PAN-2026 data with adversarially perturbed RAID examples dramatically reduces the domain gap and improves out-of-distribution robustness. Third, two-stage training is an effective strategy for fine-tuning large encoders in low-data regimes: Stage~1 produces a stable initialisation that prevents catastrophic forgetting during Stage~2 fine-tuning.

RADAR achieves a mean PAN score of 0.978 on the combined validation set and 0.997 on the PAN-2026 validation subset, exceeding the organiser-provided TF-IDF baseline by over 30 percentage points. On the RAID cross-domain test set, RADAR retains a mean of 0.972, demonstrating that the learned representations transfer well to unseen domains and attack conditions.

\textbf{Limitations and Future Work.} Several avenues remain open for improvement. Short texts ($< 100$ words) challenge both streams, and extending stylometric features with sentence-level discourse models or retrieval-augmented context could help. The paraphrase attack class remains the hardest to handle; adversarial training using a dedicated paraphrase model may close this gap. Finally, the current architecture uses a fixed 512-token window; adapting RADAR to exploit ModernBERT's full 8,192-token capacity for long-form documents is a natural extension. As LLMs continue to advance and new obfuscation strategies emerge, maintaining a live model hub with periodic re-training on fresh data will be essential to sustain detection performance over time.

